%!TEX program = uplatex
\RequirePackage{plautopatch}
\documentclass[uplatex,dvipdfmx]{jlreq}
\usepackage{amsmath,amssymb}

\pagestyle{empty}
\begin{document}

\noindent
{\large\bfseries トロイダル・コンパクト化について I, II}

\medskip
\hspace*{\fill}{\large\bfseries 佐武 一郎（中央大学・理工）}

\bigskip\bigskip

数論的商空間のコンパクト化として最初に考えられたもの（佐武・Baily--Borel のコンパクト化）は
標準的で最小という利点がある反面、非常に悪い特異点をもつものであった。
この欠点を除くために 1970 年代に井草、Hirzebruch 等により様々な非特異化が考えられたが、
その中で最も一般的かつ優れたものが、Mumford 達によって構成されたトロイダルコンパクト化である。

$\mathcal{D}$ を対称領域、$G$ をその自己同型群（の連結成分）、
$\Gamma$ をその（一つの $\mathbf{Q}$ 構造に関する）数論的部分群とするとき、
商空間 $\mathcal{V} = \Gamma \backslash \mathcal{D}$ を数論的商空間、
そのコンパクト化を数論的多様体という。
標準的コンパクト化 $\mathcal{V}^{\ast}$ は $\mathcal{V}$ に
$\mathcal{D}$ の有理境界成分 $\mathcal{F}$ から生じる低次元の数論的商空間
$\Gamma_{\mathcal{F}} \backslash \mathcal{F}$ を有限個つけ加えて得られる。
トロイダル・コンパクト化 $\bar{\mathcal{V}}^M$ は、さらに $\mathcal{V}^{\ast}$ を境界に沿って
blow up し、商空間 $\Gamma_{\mathcal{F}} \backslash \mathcal{F}$ を
その上のある種のファイバー空間で置き換えて得られるものである。
それは $\mathcal{D}$ の境界成分 $\mathcal{F}$ から生じる $\mathcal{V}$ の「穴」を埋めるために、
（$\mathcal{F}$ に関するジーゲル領域表示を使って）そのトーラス部分を
トーリック完備化で置き換えることによって構成される。

\medskip

このようにトロイダル・コンパクト化は対称領域の群論的構造をより精密に反映するものであり、
その境界上のファイバー空間の中にはアーベル多様体をファイバーとする空間（久賀のファイバー空間）、
さらにその上のトーリック束などが現れる。
したがってその構造を分析することは、これらファイバー空間およびそのコンパクト化
（いいかえれば、特異ファイバー）を調べる上に役立つと思われる。
またその上の line bundle を考察することによって、
種々の保型形式の間の関係（たとえば、Maass 形式と Jacobi 形式の関係）、
その幾何学的意味を解明できるかもしれない。
また Faltings による $\mathbf{Z}$ 上のスキームとしての構成が示すように、
数論的研究の対象としても興味あるものと思われる。

\medskip

例えていえば、この怪獣の代数幾何学的・数論的性質をもっと良く知る必要がある。
その健胆な胃の中には様々な動物達（上記のファイバー空間のコンパクト化、
そのファイバーであるアーベル多様体やトーリック束の退化したもの等）が飲み込まれている。
これら哀れな動物達を怪獣の腹中から助け出し、その本来の姿を
（より一般な枠組みの中で）回復させることは意味のある問題であろう。

\bigskip

第一の講演では、準備として、有界対称領域の境界成分、対応する放物的部分群の構造、
ジーゲル領域としての表示、典型的な例としてジーゲル空間の場合、等について概説する。

第二の講演では、数論的部分群 $\Gamma$ が与えられたとき、それに関する凸錐の分解、
トーリックな完備化、商空間 $\mathcal{V} = \Gamma \backslash \mathcal{D}$ の
$\mathcal{F}$ 方向の局所的完備化、それらの張り合わせによる全体的コンパクト化の構成等を
（主としてジーゲル空間の場合の例について）概説する。

\bigskip

この講演の補足として、三人の方々に講演をお願いすることにした。
第一日目に、トロイダル・コンパクト化の方法上また発想上、本質的なトーリック多様体について、
石井志保子さんに入門的な講演をお願いした。
第二日目にトロイダル・コンパクト化の重要な応用として、
上にも述べたアーベル多様体の退化の理論について、
その系統的研究を最初に手がけられた浪川幸彦氏に説明して頂くことにした。
最後に $\mathbf{Z}$ 上の構成については、この方面に詳しい藤原一宏氏に解説をお願いした。

\end{document}